\paragrafos{Στοιχεία και είδη τριγώνων}{Β.3.1, «Στοιχεία τριγώνου – Είδη τριγώνων»}

\begin{anagke}
Στο εργαστήριο τεχνολογίας του σχολείου η Αγγελική και ο Νίκος φτιάχνουν μια μικρή γέφυρα από ξυλάκια. Ο Νίκος πρώτα φτιάχνει ένα τετράγωνο, αλλά αυτό «κουνιέται» και στραβώνει. Η κυρία Ελένη λέει: «Βάλε ένα ξυλάκι διαγώνια και θα γίνουν δύο τρίγωνα. Το τρίγωνο δεν παραμορφώνεται!».

Η Πολυξένη μετράει τα τρίγωνα της γέφυρας και βλέπει ότι δεν είναι όλα ίδια: άλλα έχουν όλες τις πλευρές ίσες, άλλα δύο ίσες, άλλα καμία. Ο Βασίλης προσθέτει ότι άλλα έχουν ορθή γωνία και άλλα όχι. Για να συνεννοούνται, χρειάζονται ονόματα: ποια τρίγωνα λέμε ισοσκελή, ποια ορθογώνια, ποια οξυγώνια.

Σε αυτή την παράγραφο μαθαίνουμε τα στοιχεία του τριγώνου, ταξινομούμε τα τρίγωνα με δύο κριτήρια (πλευρές και γωνίες) και γνωρίζουμε τη διάμεσο, το ύψος και τη διχοτόμο.
\end{anagke}

\begin{lexi}{τρίγωνο, ισοσκελές, σκαληνό, διάμεσος}{ονόματα τριγώνων}
Το \lexiword{τρίγωνο} έχει τρεις γωνίες. Το \lexiword{ισοσκελές} έχει «ίσα σκέλη» (πλευρές), το \lexiword{ισόπλευρο} όλες τις πλευρές ίσες. Το \lexiword{σκαληνό} (από το «σκαληνός» = λοξός, ανώμαλος) δεν έχει δύο ίσες πλευρές. Η \lexiword{διάμεσος} χωρίζει την πλευρά στη «μέση».
\end{lexi}

\begin{metafrasi}
Τρεις κορυφές, τρεις πλευρές, τρεις γωνίες. & $\mathrm{Α},\mathrm{Β},\Gamma$, $\mathrm{ΑΒ},\mathrm{ΒΓ},\Gamma\mathrm{Α}$, $\hat{\mathrm{Α}},\hat{\mathrm{Β}},\hat{\Gamma}$\\
Ένα τρίγωνο με ορθή γωνία. & ορθογώνιο ($\hat{\mathrm{Α}}=90^\circ$)\\
Ένα τρίγωνο με όλες τις πλευρές ίσες. & ισόπλευρο\\
Ένα τμήμα από κορυφή στο μέσο της απέναντι πλευράς. & διάμεσος\\
Ένα τμήμα από κορυφή κάθετο στην απέναντι πλευρά. & ύψος\\
\end{metafrasi}

\tmima{Η θεωρία}

\begin{orismos}{στοιχεία τριγώνου}
Κάθε τρίγωνο $\mathrm{ΑΒΓ}$ έχει τρεις \textbf{κορυφές} $\mathrm{Α},\mathrm{Β},\Gamma$, τρεις \textbf{πλευρές} $\mathrm{ΑΒ},\mathrm{ΒΓ},\Gamma\mathrm{Α}$ και τρεις \textbf{γωνίες} $\hat{\mathrm{Α}},\hat{\mathrm{Β}},\hat{\Gamma}$. Τα $\mathrm{ΑΒ},\mathrm{ΒΓ},\Gamma\mathrm{Α}$ δηλώνουν και τα μήκη των πλευρών. Περίμετρος είναι το άθροισμα $\mathrm{ΑΒ}+\mathrm{ΒΓ}+\Gamma\mathrm{Α}$.
\end{orismos}

\begin{kanonas}{ταξινόμηση τριγώνων}
\textbf{Ως προς τις γωνίες:} \emph{ορθογώνιο} (μία ορθή γωνία), \emph{αμβλυγώνιο} (μία αμβλεία γωνία), \emph{οξυγώνιο} (όλες οι γωνίες οξείες).\\
\textbf{Ως προς τις πλευρές:} \emph{ισόπλευρο} (τρεις ίσες), \emph{ισοσκελές} (δύο ίσες), \emph{σκαληνό} (όλες άνισες).
\end{kanonas}

\begin{center}
\begin{tikzpicture}[scale=0.62]
\draw[thick] (0,0)--(0,2.2)--(2.8,0)--cycle; \node at (1,-0.6) {\small ορθογώνιο};
\begin{scope}[xshift=4.6cm]
\draw[thick] (0,0)--(3,0)--(-0.8,1.6)--cycle; \node at (1.1,-0.6) {\small αμβλυγώνιο};
\end{scope}
\begin{scope}[xshift=9.6cm]
\draw[thick] (0,0)--(3,0)--(1.4,2.4)--cycle; \node at (1.5,-0.6) {\small οξυγώνιο};
\end{scope}
\end{tikzpicture}
\end{center}

\begin{center}
\begin{tikzpicture}[scale=0.62]
\draw[thick] (0,0)--(3,0)--(1.5,2.6)--cycle; \node at (1.5,-0.6) {\small ισόπλευρο};
\begin{scope}[xshift=5cm]
\draw[thick] (0,0)--(3,0)--(1.5,2.6)--cycle; \node at (1.5,-0.6) {\small ισοσκελές};
\end{scope}
\begin{scope}[xshift=10cm]
\draw[thick] (0,0)--(3.3,0)--(0.9,2.2)--cycle; \node at (1.5,-0.6) {\small σκαληνό};
\end{scope}
\end{tikzpicture}
\end{center}

\begin{orismos}{δευτερεύοντα στοιχεία}
\textbf{Διάμεσος} τριγώνου: τμήμα που ενώνει μια κορυφή με το μέσο της απέναντι πλευράς.\\
\textbf{Ύψος}: τμήμα από μια κορυφή, κάθετο στην ευθεία της απέναντι πλευράς.\\
\textbf{Διχοτόμος}: το τμήμα της διχοτόμου μιας γωνίας του τριγώνου από την κορυφή ως την απέναντι πλευρά.
\end{orismos}

Κάθε τρίγωνο έχει τρεις διαμέσους, τρία ύψη και τρεις διχοτόμους. Οι τρεις διάμεσοι τέμνονται σε ένα σημείο (βαρύκεντρο), τα τρία ύψη επίσης (ορθόκεντρο). Στο αμβλυγώνιο τρίγωνο τα ύψη φέρονται εκτός τριγώνου, πάνω στις προεκτάσεις των πλευρών. Στο ορθογώνιο τρίγωνο οι δύο κάθετες πλευρές είναι και ύψη.

\begin{prosoxi}
\begin{itemize}[leftmargin=1.2em,itemsep=0.3ex,topsep=2pt]
\item Το ισόπλευρο είναι και ισοσκελές (έχει δύο ίσες πλευρές), αλλά όχι ανάποδα.
\item Ένα τρίγωνο έχει το πολύ μία ορθή ή μία αμβλεία γωνία.
\item Το ύψος είναι \emph{κάθετο} τμήμα, η διάμεσος πάει στο \emph{μέσο}. Δεν είναι πάντα το ίδιο τμήμα.
\item Το ύψος μπορεί να βρίσκεται έξω από το τρίγωνο.
\end{itemize}
\end{prosoxi}

\tmima{Λυμένα παραδείγματα}

\paradeigma{Η περίμετρος και το είδος}
\tixero Τρίγωνο με πλευρές $5$, $5$ και $8$ cm.\par
\tithelo Το είδος του ως προς τις πλευρές και την περίμετρο.\par
\possindeo Δύο πλευρές είναι ίσες, άρα ισοσκελές. Περίμετρος $5+5+8=18$ cm.\par
\apantisi Ισοσκελές, περίμετρος $18$ cm.\par
\elegxos Το άθροισμα των δύο μικρότερων ($5+5=10$) είναι μεγαλύτερο από την τρίτη ($8$), άρα το τρίγωνο υπάρχει.

\paradeigma{Το ύψος}
\tixero Τρίγωνο $\mathrm{ΑΒΓ}$ με $\mathrm{ΒΓ}=10$ cm και ύψος $\mathrm{ΑΔ}=4$ cm.\par
\tithelo Σχεδίασε τα ύψη όταν το τρίγωνο είναι οξυγώνιο, ορθογώνιο ή αμβλυγώνιο.\par
\possindeo Από κάθε κορυφή φέρνω κάθετο στην απέναντι πλευρά (ή την προέκτασή της, στο αμβλυγώνιο).\par
\apantisi Τα τρία ύψη περνούν από το ίδιο σημείο (ορθόκεντρο).\par
\elegxos Στο ορθογώνιο τρίγωνο το ορθόκεντρο είναι η κορυφή της ορθής γωνίας.

\paradeigma{Η διάμεσος}
\tixero Στο τρίγωνο $\mathrm{ΑΒΓ}$ η πλευρά $\mathrm{ΒΓ}=7{,}4$ cm και η $\mathrm{ΑΜ}$ είναι διάμεσος.\par
\tithelo Το μήκος $\mathrm{ΒΜ}$.\par
\possindeo Το $\mathrm{Μ}$ είναι το μέσο της $\mathrm{ΒΓ}$, άρα $\mathrm{ΒΜ}=7{,}4:2=3{,}7$ cm.\par
\apantisi $3{,}7$ cm.\par
\elegxos $\mathrm{ΒΜ}+\mathrm{ΜΓ}=3{,}7+3{,}7=7{,}4$.

\tmima{Ασκήσεις}

\zoni{Το κατάλαβα}{ονόματα και είδη}

\ex Ονόμασε τα τρίγωνα με πλευρές: α) $4,4,4$ β) $5,5,3$ γ) $3,4,5$ ως προς τις πλευρές.

\ex Βρες την περίμετρο τριγώνου με πλευρές $6{,}5$, $7$ και $8{,}5$ cm.

\ex Σχεδίασε οξυγώνιο, ορθογώνιο και αμβλυγώνιο τρίγωνο και φέρε τα ύψη τους.

\ex Στο τρίγωνο $\mathrm{ΑΒΓ}$ με $\mathrm{ΒΓ}=4{,}4$ cm φέρε τη διάμεσο $\mathrm{ΑΜ}$. Πόσο είναι το $\mathrm{ΒΜ}$;

\ex Σχεδίασε τρίγωνο $\mathrm{ΑΒΓ}$, βρες τα μέσα $\mathrm{Δ}$, $\mathrm{Ε}$, $\mathrm{Ζ}$ των πλευρών και σχεδίασε τις τρεις διαμέσους. Τι παρατηρείς;

\ex Σχεδίασε τρίγωνο και τη διχοτόμο μιας γωνίας του. Μέτρησε τις δύο γωνίες που σχηματίζονται.

\ex Σωστό ή λάθος; α) Το ισόπλευρο είναι ισοσκελές. \quad β) Το αμβλυγώνιο έχει δύο αμβλείες. \quad γ) Το ορθογώνιο τρίγωνο μπορεί να είναι ισοσκελές. \quad δ) Το ισόπλευρο είναι πάντα οξυγώνιο. \quad ε) Ένα σκαληνό δεν μπορεί να είναι ορθογώνιο.

\zoni{Το χρησιμοποιώ}{ξύλινες κατασκευές}

\ex \textbf{Η γέφυρα.} Ένα ξύλινο τρίγωνο της γέφυρας έχει πλευρές $30$, $30$ και $40$ cm. Τι είδους είναι και πόση είναι η περίμετρος;

\ex \textbf{Το μαντήλι.} Ένα τριγωνικό μαντήλι έχει πλευρές $60$, $60$ και $60$ cm. Πόσα εκατοστά κορδόνι χρειάζεται για να το γυρίσει κανείς γύρω γύρω;

\ex \textbf{Το σχοινί.} Ένα σχοινί $24$ μ. σχηματίζει ισόπλευρο τρίγωνο. Πόση είναι η πλευρά του;

\ex \textbf{Το τρίγωνο της κορνίζας.} Ισοσκελές τρίγωνο έχει περίμετρο $50$ cm και βάση $20$ cm. Πόσο μακριές είναι οι δύο ίσες πλευρές;

\ex \textbf{Η κατασκευή.} Μπορείς να φτιάξεις τρίγωνο με ξυλάκια $3$, $4$ και $9$ cm; Και με $3$, $4$ και $5$ cm; Εξήγησε.

\zoni{Το σκέφτομαι}{σύγκριση ειδών}

\ex Σε ποια είδη τριγώνων ως προς τις γωνίες μπορεί να ανήκει ένα ισοσκελές τρίγωνο;

\ex Μπορεί ένα ορθογώνιο τρίγωνο να είναι ισόπλευρο; Εξήγησε.

\ex Πόσα το πολύ ύψη του τριγώνου βρίσκονται εκτός αυτού; (Σκέψου το αμβλυγώνιο.)

\ex Συμπλήρωσε πίνακα: ποιοι συνδυασμοί (ορθογώνιο, αμβλυγώνιο, οξυγώνιο) × (σκαληνό, ισοσκελές, ισόπλευρο) υπάρχουν;

\zoni{Για το σπίτι}{τρίγωνα στην πόλη}

\ex Βρες πέντε τριγωνικά αντικείμενα στο σπίτι (σήμα, φέτα πίτσας, ράφι, στέγη κ.λπ.) και πες το είδος του καθενός.

\ex Κόψε από χαρτί τρίγωνο και δίπλωσε το ώστε να βρεις το μέσο μιας πλευράς και τη διάμεσο.
